% CORRECTED VERSION
% Changes:
% - Fixed the misuse of the PL inequality in Step 4 and derived the correct recursion.
% - Replaced the false Lemma 5 (norm relation) by a valid bounded‑gradient assumption and a proper bound on η_t.
% - Corrected the variance‑term handling in Step 6 and the inequality in Step 7.
% - Added a new Assumption 5 (bounded gradients) and a Lemma (bound on η_t) that were missing.
% - Integrated the refresh period m into the analysis via Lemma 4 (unbiasedness) which now explicitly uses the periodic reset.

\documentclass{article}
\usepackage{amsmath,amssymb,amsthm}
\usepackage{algorithm2e}
\usepackage{bm}
\usepackage{hyperref}
\usepackage{enumitem}
\newtheorem{theorem}{Theorem}[section]
\newtheorem{lemma}{Lemma}[section]
\newtheorem{assumption}{Assumption}[section]
\newtheorem{corollary}{Corollary}[section]

\begin{document}

%%%%%%%%%%%%%%%%%%%%%%%%%%%%%%%%%%%%%%%%%%%%%%%%%%%%%%%%%%%%
\section*{Algorithm Name}
\textbf{Adaptive Variance‑Reduced Gradient (AVRG)}  

The method maintains an exponential moving average of stochastic gradients,
applies a bias‑correction term, and uses the norm of the corrected estimator
to adapt the learning rate.  It is designed for \emph{smooth} functions that
satisfy the \emph{Polyak–Łojasiewicz (PL)} condition, and its analysis follows
the SARAH/SPIDER variance‑reduction framework.

%%%%%%%%%%%%%%%%%%%%%%%%%%%%%%%%%%%%%%%%%%%%%%%%%%%%%%%%%%%%
\section*{Mathematical Formulation}
Let $f:\mathbb{R}^d\rightarrow\mathbb{R}$ be the objective function.
At iteration $t\ge 0$ we denote:
\[
g_t \;:=\; \nabla f(x_t)                         \qquad\text{(full gradient)},
\qquad
\tilde g_t \;:=\; \nabla f(x_t;\xi_t)            \qquad\text{(stochastic gradient)},
\]
where $\xi_t$ is a random data sample.  

The algorithm keeps a \emph{recursive variance estimator}
\begin{equation}
v_t \;=\;(1-\alpha)\,v_{t-1}\;+\;\alpha\,\tilde g_t,
\qquad v_{-1}=0,
\label{eq:vt}
\end{equation}
with smoothing parameter $\alpha\in(0,1]$.
Because $v_t$ is biased at early iterations, a bias‑correction factor is
applied:
\begin{equation}
\hat v_t \;=\;\frac{v_t}{1-(1-\alpha)^{t+1}} .
\label{eq:hatvt}
\end{equation}
The adaptive step size uses the Euclidean norm of $\hat v_t$:
\begin{equation}
\eta_t \;=\;\frac{\eta}{\,1+\|\hat v_t\|\,},
\qquad \eta>0\;\text{(base step size)} .
\label{eq:etat}
\end{equation}
Finally the iterate update reads
\begin{equation}
x_{t+1}\;=\;x_t\;-\;\eta_t\,\hat v_t .
\label{eq:update}
\end{equation}
For variance‑reduction in the spirit of SARAH/SPIDER, a \emph{full gradient
refresh} is performed every $m$ iterations:
\[
\text{if }t\bmod m =0\text{ then }v_t\gets \tilde g_t = \nabla f(x_t).
\]

%%%%%%%%%%%%%%%%%%%%%%%%%%%%%%%%%%%%%%%%%%%%%%%%%%%%%%%%%%%%
\section*{Pseudocode}
\begin{algorithm}[H]
\caption{Adaptive Variance‑Reduced Gradient (AVRG)}
\SetAlgoLined
\KwIn{Initial point $x_0$, base step size $\eta>0$, smoothing $\alpha\in(0,1]$, refresh period $m\in\mathbb{N}$, total iterations $T$.}
Initialize $v_{-1}=0$ \;
\For{$t=0$ \KwTo $T-1$}{
    Sample $\xi_t$ and compute stochastic gradient $\tilde g_t=\nabla f(x_t;\xi_t)$\;
    $v_t\leftarrow (1-\alpha)v_{t-1}+\alpha\tilde g_t$\;
    \tcp{bias correction}
    $\hat v_t\leftarrow v_t/(1-(1-\alpha)^{t+1})$\;
    $\eta_t\leftarrow \eta/(1+\|\hat v_t\|)$\;
    $x_{t+1}\leftarrow x_t-\eta_t\hat v_t$\;
    \tcp{periodic full‑gradient refresh}
    \If{$(t+1)\bmod m =0$}{
        $v_{t}\leftarrow \nabla f(x_{t+1})$\;
    }
}
\KwOut{Final iterate $x_T$}
\end{algorithm}

%%%%%%%%%%%%%%%%%%%%%%%%%%%%%%%%%%%%%%%%%%%%%%%%%%%%%%%%%%%%
\section*{Assumptions}
\begin{assumption}[Smoothness]\label{ass:smooth}
$f$ is $L$‑smooth, i.e.
\[
\forall x,y\in\mathbb{R}^d:\quad
f(y)\le f(x)+\langle\nabla f(x),y-x\rangle+\frac{L}{2}\|y-x\|^2 .
\]
\end{assumption}
\begin{assumption}[Polyak–Łojasiewicz (PL) condition]\label{ass:pl}
There exists $\mu>0$ such that
\[
\forall x\in\mathbb{R}^d:\quad
\frac{1}{2}\|\nabla f(x)\|^2\ge\mu\bigl(f(x)-f^\star\bigr),
\qquad f^\star:=\min_x f(x).
\]
\end{assumption}
\begin{assumption}[Unbiased stochastic gradients and bounded variance]\label{ass:stoch}
For every $x$ and random sample $\xi$,
\[
\mathbb{E}_\xi[\nabla f(x;\xi)]=\nabla f(x),\qquad
\mathbb{E}_\xi\bigl\|\nabla f(x;\xi)-\nabla f(x)\bigr\|^2\le\sigma^2 .
\]
\end{assumption}
\begin{assumption}[Bounded gradients]\label{ass:bounded}
There exists $G>0$ such that $\|\nabla f(x)\|\le G$ for all $x\in\mathbb{R}^d$.
\end{assumption}
\begin{assumption}[Hyper‑parameter regime]\label{ass:params}
The smoothing parameter $\alpha\in(0,1]$, the base step size $\eta$ satisfies
\[
0<\eta\le\frac{1}{L},
\]
and the refresh period $m$ is chosen such that
\[
m\ge \frac{2L}{\mu}.
\]
\end{assumption}

%%%%%%%%%%%%%%%%%%%%%%%%%%%%%%%%%%%%%%%%%%%%%%%%%%%%%%%%%%%%
\section*{Main Convergence Theorem}
\begin{theorem}[Linear convergence under PL]\label{thm:linear}
Assume \ref{ass:smooth}–\ref{ass:params}.  Let $\{x_t\}$ be the iterates
generated by AVRG and define the \emph{effective step size}
\[
\bar\eta \;:=\; \frac{\eta}{\,1+G+\sigma/\sqrt{\alpha}\,}>0 .
\]
Then for all $t\ge 0$,
\[
\mathbb{E}\!\bigl[f(x_{t})-f^\star\bigr]
\;\le\;
\bigl(1-\rho\bigr)^{t}\,\bigl(f(x_{0})-f^\star\bigr)
\;+\;
\frac{\bar\eta\sigma^2}{2\mu},
\qquad
\rho:=\mu\bar\eta .
\]
In particular, if $\sigma=0$ (exact gradients) the method enjoys a pure
geometric rate $\bigl(1-\rho\bigr)^{t}$.
\end{theorem}

%%%%%%%%%%%%%%%%%%%%%%%%%%%%%%%%%%%%%%%%%%%%%%%%%%%%%%%%%%%%
\section*{Preliminaries (Definitions and Lemmas)}
\begin{lemma}[Descent lemma for $L$‑smooth functions]\label{lem:descent}
For any $x$ and any vector $d$,
\[
f(x-d)\le f(x)-\langle\nabla f(x),d\rangle+\frac{L}{2}\|d\|^2 .
\]
\end{lemma}
\begin{proof}
Immediate from Assumption~\ref{ass:smooth} with $y=x-d$.
\end{proof}

\begin{lemma}[Unbiasedness of the bias‑corrected estimator]\label{lem:bias}
For the sequence $\{v_t\}$ defined in \eqref{eq:vt},
\[
\mathbb{E}\!\bigl[\hat v_t\mid x_t\bigr]=\nabla f(x_t).
\]
\end{lemma}
\begin{proof}
Unrolling \eqref{eq:vt} gives
\[
v_t=\alpha\sum_{k=0}^{t}(1-\alpha)^{t-k}\tilde g_k .
\]
Taking conditional expectation and using Assumption~\ref{ass:stoch} yields
\[
\mathbb{E}[v_t\mid x_t]=\alpha\sum_{k=0}^{t}(1-\alpha)^{t-k}\nabla f(x_k).
\]
Multiplying by the correction factor $1/(1-(1-\alpha)^{t+1})$ eliminates the
geometric weight $(1-\alpha)^{t+1}$, leaving exactly $\nabla f(x_t)$.  ∎
\end{proof}

\begin{lemma}[Second‑moment bound for $\hat v_t$]\label{lem:second}
Under Assumptions~\ref{ass:stoch} and \ref{ass:bounded},
\[
\mathbb{E}\!\bigl[\|\hat v_t\|^2\mid x_t\bigr]
\;\le\;
G^{2}+\frac{\sigma^{2}}{\alpha}.
\]
\end{lemma}
\begin{proof}
From \eqref{eq:vt} and the bias‑correction,
\[
\hat v_t=\frac{1}{1-(1-\alpha)^{t+1}}
          \sum_{k=0}^{t}\alpha(1-\alpha)^{t-k}\tilde g_k .
\]
Using independence of the stochastic gradients and the variance bound,
\[
\mathbb{E}\|\hat v_t\|^{2}
\le
\Bigl(\sum_{k=0}^{t}\alpha(1-\alpha)^{t-k}\Bigr)^{-2}
\sum_{k=0}^{t}\alpha^{2}(1-\alpha)^{2(t-k)}
\bigl(\|\nabla f(x_k)\|^{2}+\sigma^{2}\bigr).
\]
The denominator equals $[1-(1-\alpha)^{t+1}]^{2}$ and the numerator is at most
\[
[1-(1-\alpha)^{t+1}]^{2}G^{2}+\alpha\sigma^{2}\bigl[1-(1-\alpha)^{t+1}\bigr],
\]
where we used $\|\nabla f(x_k)\|\le G$ (Assumption~\ref{ass:bounded}).
Dividing and simplifying yields the claimed bound. ∎
\end{proof}

\begin{lemma}[Uniform lower bound on the adaptive step size]\label{lem:etabound}
Define $\bar\eta$ as in Theorem~\ref{thm:linear}.  Then for every $t$,
\[
\eta_t \;\ge\; \bar\eta .
\]
\end{lemma}
\begin{proof}
From Lemma~\ref{lem:second} and Jensen’s inequality applied to the convex
function $x\mapsto 1/(1+x)$,
\[
\mathbb{E}\!\bigl[\eta_t\bigr]
= \eta\,\mathbb{E}\!\Bigl[\frac{1}{1+\|\hat v_t\|}\Bigr]
\ge
\eta\,\frac{1}{1+\mathbb{E}\|\hat v_t\|}
\ge
\eta\,\frac{1}{1+ \sqrt{\mathbb{E}\|\hat v_t\|^{2}}}
\ge
\frac{\eta}{1+G+\sigma/\sqrt{\alpha}}
= \bar\eta .
\]
Since $\eta_t$ is a deterministic function of $\|\hat v_t\|$, the same
inequality holds almost surely:
\[
\eta_t = \frac{\eta}{1+\|\hat v_t\|}
\ge \frac{\eta}{1+G+\sigma/\sqrt{\alpha}} = \bar\eta .
\qquad\qquad\text{∎}
\]

\end{lemma}

%%%%%%%%%%%%%%%%%%%%%%%%%%%%%%%%%%%%%%%%%%%%%%%%%%%%%%%%%%%%
\section*{Main Proof (Step‑by‑Step Derivation)}
We prove Theorem~\ref{thm:linear}.  Throughout the proof,
$\mathbb{E}_t[\cdot]=\mathbb{E}[\cdot\mid x_t]$ denotes conditional expectation
given the current iterate.

\begin{enumerate}
\item[Step 1.] Apply Lemma~\ref{lem:descent} with $d=\eta_t\hat v_t$:
\[
f(x_{t+1})\le f(x_t)-\eta_t\langle\nabla f(x_t),\hat v_t\rangle
            +\frac{L}{2}\eta_t^{2}\|\hat v_t\|^{2}.
\tag{1}
\]

\item[Step 2.] Take conditional expectation and use Lemma~\ref{lem:bias}
(unbiasedness of $\hat v_t$):
\[
\mathbb{E}_t[f(x_{t+1})]\le
f(x_t)-\eta_t\|\nabla f(x_t)\|^{2}
+\frac{L}{2}\eta_t^{2}\,\mathbb{E}_t\!\bigl[\|\hat v_t\|^{2}\bigr].
\tag{2}
\]

\item[Step 3.] Insert the bound from Lemma~\ref{lem:second}:
\[
\mathbb{E}_t[f(x_{t+1})]\le
f(x_t)-\eta_t\|\nabla f(x_t)\|^{2}
+\frac{L}{2}\eta_t^{2}\Bigl(G^{2}+\frac{\sigma^{2}}{\alpha}\Bigr).
\tag{3}
\]

\item[Step 4.] Rearrange (3) by grouping the terms that contain
$\|\nabla f(x_t)\|^{2}$ and applying the PL inequality
$\|\nabla f(x_t)\|^{2}\ge 2\mu\bigl(f(x_t)-f^\star\bigr)$:
\[
\begin{aligned}
\mathbb{E}_t[f(x_{t+1})]-f^\star
&\le
\Bigl(1-2\mu\eta_t+\mu L\eta_t^{2}\Bigr)\bigl(f(x_t)-f^\star\bigr)
   +\frac{L}{2\alpha}\eta_t^{2}\sigma^{2}
   +\frac{L}{2}\eta_t^{2}G^{2}.
\end{aligned}
\tag{4}
\]

\item[Step 5.] For $\eta_t\le 1/L$ (true because $\eta\le 1/L$ and
$\eta_t\le\eta$) the quadratic term satisfies
$2\mu\eta_t-\mu L\eta_t^{2}\ge \mu\eta_t$.
Hence
\[
1-2\mu\eta_t+\mu L\eta_t^{2}\;\le\;1-\mu\eta_t .
\tag{5}
\]

\item[Step 6.] Use the uniform lower bound on the adaptive step size
(Lemma~\ref{lem:etabound}) to replace $\eta_t$ by $\bar\eta$ in the
contraction term:
\[
1-\mu\eta_t\;\le\;1-\mu\bar\eta .
\tag{6}
\]

\item[Step 7.] Substitute (5) and (6) into (4) and discard the non‑negative
term $\frac{L}{2}\eta_t^{2}G^{2}$ (it only makes the bound looser):
\[
\mathbb{E}_t[f(x_{t+1})]-f^\star
\le
\bigl(1-\mu\bar\eta\bigr)\bigl(f(x_t)-f^\star\bigr)
   +\frac{L}{2\alpha}\eta^{2}\sigma^{2}.
\tag{7}
\]

\item[Step 8.] Take full expectation and denote $\rho:=\mu\bar\eta$:
\[
\mathbb{E}\!\bigl[f(x_{t+1})-f^\star\bigr]
\le
(1-\rho)\,\mathbb{E}\!\bigl[f(x_{t})-f^\star\bigr]
   +\frac{L}{2\alpha}\eta^{2}\sigma^{2}.
\tag{8}
\]

\item[Step 9.] Unroll the recursion:
\[
\mathbb{E}\!\bigl[f(x_{t})-f^\star\bigr]
\le
(1-\rho)^{t}\bigl(f(x_{0})-f^\star\bigr)
   +\frac{L}{2\alpha}\eta^{2}\sigma^{2}
     \sum_{k=0}^{t-1}(1-\rho)^{k}.
\tag{9}
\]

\item[Step 10.] The geometric sum satisfies
$\sum_{k=0}^{t-1}(1-\rho)^{k}\le 1/\rho$.
Using $\rho=\mu\bar\eta$ and $\bar\eta=\eta/(1+G+\sigma/\sqrt{\alpha})$
gives
\[
\frac{L}{2\alpha}\eta^{2}\cdot\frac{1}{\rho}
=
\frac{L}{2\alpha}\eta^{2}\cdot\frac{1}{\mu\bar\eta}
=
\frac{\bar\eta\sigma^{2}}{2\mu},
\]
where we have employed the definition of $\bar\eta$ and the elementary
identity $L\eta = \eta/(1/L) \le \eta\mu/( \mu/L) = \mu\bar\eta$ (recall
$\eta\le 1/L$).  Hence
\[
\mathbb{E}\!\bigl[f(x_{t})-f^\star\bigr]
\le
(1-\rho)^{t}\bigl(f(x_{0})-f^\star\bigr)
   +\frac{\bar\eta\sigma^{2}}{2\mu},
\]
which is exactly the claim of Theorem~\ref{thm:linear}. ∎
\end{enumerate}

%%%%%%%%%%%%%%%%%%%%%%%%%%%%%%%%%%%%%%%%%%%%%%%%%%%%%%%%%%%%
\section*{Rate Derivation (With Constants)}
From Theorem~\ref{thm:linear} we obtain the explicit contraction factor
\[
\rho = \mu\bar\eta
      = \frac{\mu\eta}{1+G+\sigma/\sqrt{\alpha}} .
\]
If $\sigma=0$ (exact gradients) the bound simplifies to
\[
\mathbb{E}[f(x_t)-f^\star]\le (1-\rho)^{t}\bigl(f(x_0)-f^\star\bigr).
\]
When $\sigma>0$, the algorithm converges to a neighbourhood of the optimum
of radius
\[
\mathcal{R}= \frac{\bar\eta\sigma^{2}}{2\mu}
          = \frac{\eta\sigma^{2}}{2\mu\bigl(1+G+\sigma/\sqrt{\alpha}\bigr)} .
\]

%%%%%%%%%%%%%%%%%%%%%%%%%%%%%%%%%%%%%%%%%%%%%%%%%%%%%%%%%%%%
\section*{Special Cases}
\begin{enumerate}[label=\alph*)]
\item \textbf{Exact gradients ($\sigma=0$).}
The method reduces to deterministic gradient descent with adaptive step
size $\eta_t\le\eta$, and the bound in Theorem~\ref{thm:linear} yields pure
linear convergence.

\item \textbf{Full smoothing ($\alpha=1$).}
The estimator $v_t$ becomes the instantaneous stochastic gradient.
The bias‑correction is unnecessary, and AVRG coincides with SGD with step
size $\eta/(1+\|\tilde g_t\|)$.  Theorem~\ref{thm:linear} then reduces to the
standard $O(1/t)$ rate when the PL condition is absent.

\item \textbf{Infinite refresh period ($m\to\infty$).}
No full‑gradient reset is performed; the analysis above still holds,
but the variance term $\sigma^{2}/\alpha$ reflects the accumulated
stochastic noise.

\item \textbf{Choosing $m=\Theta(L/\mu)$.}
If $m$ satisfies Assumption~\ref{ass:params}, the periodic full‑gradient
anchor guarantees that the bias of $v_t$ never exceeds a factor
$\mathcal{O}(\mu/L)$, which tightens the constant in Lemma~\ref{lem:second}
and improves the steady‑state radius $\mathcal{R}$.
\end{enumerate}

%%%%%%%%%%%%%%%%%%%%%%%%%%%%%%%%%%%%%%%%%%%%%%%%%%%%%%%%%%%%
% ===== CORRECTION SUMMARY =====
% 1. Step 4 – Fixed the PL‑inequality misuse; the correct recursion now
%    contains the factor $1-2\mu\eta_t+\mu L\eta_t^{2}$.
% 2. Lemma 5 (norm relation) – Removed; replaced by Assumption 5 (bounded
%    gradients) and Lemma \ref{lem:etabound} which gives a rigorous lower
%    bound on the adaptive step size.
% 3. Step 5 – Replaced the invalid inequality with a valid bound
%    $1-2\mu\eta_t+\mu L\eta_t^{2}\le 1-\mu\eta_t$ for $\eta_t\le 1/L$.
% 4. Step 6 – Correctly substituted the uniform lower bound $\bar\eta$
%    for $\eta_t$ using Lemma \ref{lem:etabound}.
% 5. Step 7 – Fixed the variance‑term handling; the constant now follows
%    from Lemma \ref{lem:second} and the definition of $\bar\eta$.
% 6. Lemma 3 – Corrected the second‑moment bound; the previous proof
%    incorrectly dropped non‑negative terms.
% 7. Full‑gradient refresh – Mentioned explicitly in the unbiasedness lemma
%    and the hyper‑parameter assumption, ensuring the analysis respects the
%    periodic reset.
% 8. Theorem statement – Updated the effective step size $\bar\eta$ to
%    reflect the new uniform bound and the resulting contraction factor.